\documentclass{article}
\usepackage[utf8]{inputenc}
\usepackage[T1]{fontenc}
\usepackage{geometry}
\usepackage{graphicx}
\usepackage{float}
\usepackage{hyperref}
\usepackage[english,indonesian]{babel}

\title{Optimasi Pemrosesan Template LaTeX Menggunakan Pendekatan Berbasis AI}
\author{Academic Support Agent}
\date{\today}

\begin{document}
\maketitle

\begin{abstract}

\end{abstract}

\chapter{PENDAHULUAN}


\chapter{TINJAUAN PUSTAKA}


% Original content (commented out):
% \documentclass[10pt,twocolumn,letterpaper]{article}
%% Welcome to Overleaf!
%% If this is your first time using LaTeX, it might be worth going through this brief presentation:
%% https://www.overleaf.com/latex/learn/free-online-introduction-to-latex-part-1

%% Researchers have been using LaTeX for decades to typeset their papers, producing beautiful, crisp documents in the process. By learning LaTeX, you are effectively following in their footsteps, and learning a highly valuable skill!

%% The \usepackage commands below can be thought of as analogous to importing libraries into Python, for instance. We've pre-formatted this for you, so you can skip right ahead to the title below.

%% Language and font encodings
\usepackage[spanish,english]{babel}
\usepackage[utf8x]{inputenc}
\usepackage[T1]{fontenc}

%% Sets page size and margins
\usepackage[a4paper,top=3cm,bottom=2cm,left=3cm,right=3cm,marginparwidth=1.75cm]{geometry}

%% Useful packages
\usepackage{amsmath}
\usepackage{graphicx}
\usepackage[colorinlistoftodos]{todonotes}
\usepackage[colorlinks=true, allcolors=blue]{hyperref}

%% Title
\title{
		%\vspace{-1in} 	
		\usefont{OT1}{bch}{b}{n}
		\normalfont \normalsize \textsc{STEM Fellowship Journal Template} \\ [14pt]
		\huge Title goes here! \\
}

\usepackage{authblk}

\author[1]{First Author}
\author[2]{Second Author}
\author[1]{Third Author}
\author[1,3]{Fourth Author}

	\affil[1]{\small{Department, Institution/School}}
\affil[2]{Department, Institution/School}
\affil[3]{Department, Institution/School} 

\begin{document}
\maketitle

\selectlanguage{english}
\begin{abstract}
The abstract is in a sense the most important part of your paper, since many readers will scan an abstract first before deciding to invest effort in reading the paper.

An abstract should be about 300 words long at most, and should act as a clear summary of the paper. It should state the aim and scope of the research, methods, results and conclusions, and the implications of the paper’s finding. The abstract should be broadly accessible (i.e. able to be understood by as many people as possible - even those outside the field) and communicate the importance of the work being done. 

Structure the abstract as follows: background and aims, methods, results and conclusion. Please note that review articles will not contain methods or results. 

Examples of ideal abstracts can be found \href{http://journal.stemfellowship.org/userimages/ContentEditor/1451371646443/AbstractExamples.pdf}{here}, with the main ideas highlighted.
\end{abstract} 

\section*{Keywords}
Select keywords with care, because they will help users discover your paper. To determine appropriate keywords, put yourself in the position of someone who is trying to search for a paper like yours. What search terms would you use? From those terms, select a list of at least three and no more than five words. Include these words in the text of the abstract and if at all possible, in the title of the paper.

\section{Introduction}
Before we get to the actual introduction, welcome to Overleaf, as well as \LaTeX\ itself! Although \LaTeX\ certainly has its quirks, we hope that by contrasting the template you see here with the compiled document on the right side, you can get an intuitive sense of how to work with it.

Another thing before the introduction; here, I'm going attach a citation to this sentence \cite{a_reference}. Scroll on down to the bibliography section of the \LaTeX\ code if you'd like to see the other end of the built-in references system. The numbering is all handled in-house -- you just have to assign each reference a key, and Overleaf takes care of the rest.

On with the actual introduction! The introduction should provide context and background information appropriate for an academic audience. It should state a focused research question, define the variables being studied, and make clear the objective, importance, and relevance of the work.

Here is where you'd introduce the context surrounding your study. What led you to the question you ended up asking? Why is it relevant? Which fields of science is your question based around? What has previous literature demonstrated?

While the structure of the previous parts of the introduction can be relatively variable, you must make sure to provide a brief overview of the study itself, and the methods you used to accomplish it. Obviously, excessive detail is not necessary (that's what the next section is for). Lastly, be sure to make mention of the potential implications of your findings, but once again remember that you'll be going into more detail about that in the discussion.

\section{Materials \& Methods}
In this section, materials and methods should be reported in a way that they can be reproduced by other investigators. It should explain the methods used and the reasons for using them. Provide sources for materials used wherever possible.

The method should be written as a step by step process, and explain how variables will be measured. It should also state and explain the choices of calculations and techniques of statistical analysis.

Be as concise and to-the-point as possible, and remember - \textbf{do not justify your methods here!} You simply need to state what you did. You can (and probably should) mention the purpose of using a certain computational tool within the context of what you set out to achieve, but mentioning things like "it's particularly efficient at this and better than all competing computational tools" is unnecessary in the methods section. However, you can definitely talk about all of this in the discussion, and talk about why your methods are, say, the most effective ones for the task.

Think of this section as a technical manual of sorts, that another team of researchers could read and easily follow in order to replicate what you did to carry out this study.

\section{Results}

Describe the results of the experiments conducted, and include any tables or figures that report your results. Make sure to not analyze any of your results here! All of that belongs in the discussion.

It is important that you \textbf{do not fabricate or plagiarise data, nor exclude any data that is contrary to your hypothesis}. Report all observations made during the research.

Including figures into \LaTeX\ can seem intimidating at first, but Overleaf makes it easy: simply click the 'Project' button above, select 'Files', and upload away from your computer. Then, insert the file name into the appropriate section of the code below.  Figure 1  shows the output of such code. A guide to formatting figures can be found at \url{https://en.wikibooks.org/wiki/LaTeX/Floats,_Figures_and_Captions#Figures}.
\\

{\scriptsize
\begin{verbatim}
\begin{figure}
    \centering
    \includegraphics[width=0.4\textwidth]{test.png}
    \caption{Hello!}
\end{figure}
\end{verbatim}
}

\begin{figure}
  \centering
  \includegraphics[width=0.4\textwidth]{test.png}
  \caption{Notice how \LaTeX\ automatically numbers this figure.}
\end{figure}


\section{Discussion}
This is where you interpret your results, pointing out interesting trends within your data and how they relate to your initial hypothesis. This is also the place to justify your methodology, if you're so inclined (i.e. Why did you specifically use a certain statistical test over another? Why this tool over that tool?). Lastly, you're going to want to discuss potential sources of error. Make sure to make explicit reference to figures/tables when discussing your data; it can be helpful to walk the reader through your own personal interpretation of each figure in order.

Explain observations, patterns and trends as represented in graphs using scientific theory and reasoning. The discussion of your results should comprise an interpretation of the results in terms of their significance in relation to your original objective.

\section{Conclusions}
Use this section to briefly summarize findings, comment on the significance of these findings, and discuss potential future directions for study. This section should always be included in review articles, but is optional in original research articles, since conclusions — without comments — on the significance of findings or suggestions for future work can be included at the end of the discussion section. However, it is strongly recommended.

What are the long-term implications of your findings? Wrap up your discussion succinctly while pointing out the significance of your work as well as it what it means for the fields you examined as much as possible. Lastly, suggest ideas for future studies that could build on your work, and justify why they might be useful. Otherwise, you're all done!

\section{Acknowledgements}
In the Acknowledgements section, the author(s) acknowledge or thank any persons or institutions who helped support the work in any way. In particular, this section must disclose any funding for the research completed, including the grant number (if a grant was awarded).

Anyone to thank/credit for helping your team along the way? This is the place to do it.

\section{Appendix}
Appendices are optional sections that should include additional tables, figures, or other data beyond what is included in the results section. An appendix should be able to stand separately from the prinicpal article. Therefore, no references (callouts) to appendix figures or tables should be made in the principal article. If required, the appendix should have its own bibliographic reference list with citations to that list confined to the appendix. \cite{latexcompanion}

\newpage

\bibliographystyle{vancouver}
\bibliography{references}

\begin{thebibliography}{9}
All information or data from other sources or authors must be appropriately cited at each point in the text where the source is relevant. All sources should be listed in a bibliographic reference section at the end of the paper with the heading "Reference List". 

If the paper contains any copyright material, such as tables, figures, images, or other, the authors must obtain written permission from the copyright holder and provide it to the editor. 

Bibliographic references should conform to the Vancouver style. In this style, bibliographic references are numbered and citations in the text to each reference should be numbers in square brackets; e.g., [1]. 

If you are using a reference management system (such as Zotero, Mendeley or ReadCube), you may be able to export your bibliography as a Bib\TeX{} file. 

Please refer to the comments in this template for how to format your references in LaTeX (either with a .bib file, or via manual input of references).


%% Please read the following comments for instructions on formatting your references.
%% The provided style file formats references in the required Vancouver style
%% This is a useful guide to formatting bibliographies using LaTeX (https://en.wikibooks.org/wiki/LaTeX/Bibliography_Management)

%% References with BibTeX database:
%%Bib\TeX{} is a bibliographic tool that is used with \LaTeX{} to help organize the user's references and create a bibliography. This short tutorial explains how to use BibTex files [https://www.latex-tutorial.com/tutorials/beginners/latex-bibtex/].

%%For references without a BibTeX database:
%%Your line of code will start with the \bibitem{cite_key} command. This is a unique identifier for a reference, such as #\bibitem{Lindquist95} or  #\bibitem{Smithetal2002}. Following the #\bibitem{cite_key} command is the actual reference. Here you should include all relevant details for the reference e.g. authors, title, year of publication, etc. The \bibliographystyle{vancouver} will format these references to the correct Vancouver style - as required by the STEM Fellowship Journal.

%% \bibitem must have the following form:
%%   \bibitem{key}...
%%

% \bibitem{}
% Below is an example.

%%\bibitem{Smithetal2002}

\end{thebibliography}

\end{document}

\end{document}